\begin{ex}
Ki-lô-gam được viết tắt bằng mấy chữ cái trong bảng chữ cái tiếng Anh? (viết tắt là kg).
\shortans{2}
\loigiai{
Viết tắt là kg, gồm 2 chữ cái k và g.
}
\end{ex}

\begin{ex}
Tính: $5\text{ kg} + 3\text{ kg}$. Kết quả là bao nhiêu ki-lô-gam?
\shortans{8}
\loigiai{
Ta có: $5 + 3 = 8$ kg.
}
\end{ex}

\begin{ex}
Tính: $14\text{ kg} + 5\text{ kg}$. Kết quả là bao nhiêu ki-lô-gam?
\shortans{19}
\loigiai{
Ta có: $14 + 5 = 19$ kg.
}
\end{ex}

\begin{ex}
Tính: $26\text{ kg} + 12\text{ kg}$. Kết quả là bao nhiêu ki-lô-gam?
\shortans{38}
\loigiai{
Ta có: $26 + 12 = 38$ kg.
}
\end{ex}

\begin{ex}
Tính: $45\text{ kg} - 20\text{ kg}$. Kết quả là bao nhiêu ki-lô-gam?
\shortans{25}
\loigiai{
Ta có: $45 - 20 = 25$ kg.
}
\end{ex}

\begin{ex}
Tính: $18\text{ kg} - 9\text{ kg}$. Kết quả là bao nhiêu ki-lô-gam?
\shortans{9}
\loigiai{
Ta có: $18 - 9 = 9$ kg.
}
\end{ex}

\begin{ex}
Tính: $35\text{ kg} - 15\text{ kg}$. Kết quả là bao nhiêu ki-lô-gam?
\shortans{20}
\loigiai{
Ta có: $35 - 15 = 20$ kg.
}
\end{ex}

\begin{ex}
Tính giá trị biểu thức: $12\text{ kg} + 8\text{ kg} - 5\text{ kg}$.
\shortans{15}
\loigiai{
Ta có $12 + 8 = 20$, rồi $20 - 5 = 15$ kg.
}
\end{ex}

\begin{ex}
Tính giá trị biểu thức: $25\text{ kg} - 10\text{ kg} + 7\text{ kg}$.
\shortans{22}
\loigiai{
Ta có $25 - 10 = 15$, rồi $15 + 7 = 22$ kg.
}
\end{ex}

\begin{ex}
Trên đĩa cân thăng bằng: đĩa bên trái có quả dưa hấu, đĩa bên phải có 2 quả cân loại 2 kg và 1 quả cân loại 1 kg. Quả dưa hấu nặng bao nhiêu ki-lô-gam?
\shortans{5}
\loigiai{
Quả dưa nặng: $2 + 2 + 1 = 5$ kg.
}
\end{ex}

\begin{ex}
Đĩa bên trái có một túi bột mì và quả cân 1 kg, đĩa bên phải có quả cân 5 kg. Hai đĩa cân thăng bằng. Túi bột mì nặng bao nhiêu ki-lô-gam?
\shortans{4}
\loigiai{
Túi bột mì nặng: $5 - 1 = 4$ kg.
}
\end{ex}

\begin{ex}
Bao gạo to nặng 25 kg, bao gạo nhỏ nặng 15 kg. Cả hai bao gạo nặng bao nhiêu ki-lô-gam?
\shortans{40}
\loigiai{
Cả hai bao nặng: $25 + 15 = 40$ kg.
}
\end{ex}

\begin{ex}
Bao gạo to nặng hơn bao gạo nhỏ bao nhiêu ki-lô-gam?
\shortans{10}
\loigiai{
Nặng hơn: $25 - 15 = 10$ kg.
}
\end{ex}

\begin{ex}
Con ngỗng nặng 6 kg, con vịt nặng 2 kg. Con ngỗng nặng hơn con vịt bao nhiêu ki-lô-gam?
\shortans{4}
\loigiai{
Nặng hơn: $6 - 2 = 4$ kg.
}
\end{ex}

\begin{ex}
Con ngỗng và con vịt nặng tất cả bao nhiêu ki-lô-gam?
\shortans{8}
\loigiai{
Cả hai con nặng: $6 + 2 = 8$ kg.
}
\end{ex}

\begin{ex}
Con chó nặng 12 kg, con mèo nặng 4 kg. Con chó nặng hơn con mèo bao nhiêu ki-lô-gam?
\shortans{8}
\loigiai{
Con chó nặng hơn: $12 - 4 = 8$ kg.
}
\end{ex}

\begin{ex}
Mẹ mua về 18 kg gạo, gia đình đã ăn hết 9 kg gạo. Còn lại bao nhiêu ki-lô-gam gạo?
\shortans{9}
\loigiai{
Còn lại: $18 - 9 = 9$ kg.
}
\end{ex}

\begin{ex}
Một thùng cam nặng 32 kg, người ta lấy ra 10 kg cam. Thùng cam còn lại nặng bao nhiêu ki-lô-gam?
\shortans{22}
\loigiai{
Còn lại: $32 - 10 = 22$ kg.
}
\end{ex}

\begin{ex}
Điền số thích hợp: $17\text{ kg} + \dots\text{ kg} = 25\text{ kg}$.
\shortans{8}
\loigiai{
Số cần điền là: $25 - 17 = 8$.
}
\end{ex}

\begin{ex}
Điền số thích hợp: $48\text{ kg} - \dots\text{ kg} = 23\text{ kg}$.
\shortans{25}
\loigiai{
Số cần điền là: $48 - 23 = 25$.
}
\end{ex}

\begin{ex}
Điền số thích hợp: $\dots\text{ kg} + 15\text{ kg} = 39\text{ kg}$.
\shortans{24}
\loigiai{
Số cần điền là: $39 - 15 = 24$.
}
\end{ex}

\begin{ex}
Bé Nam cân nặng 23 kg, bé Mai cân nặng 21 kg. Cả hai bạn cân nặng tất cả bao nhiêu ki-lô-gam?
\shortans{44}
\loigiai{
Cả hai bạn nặng: $23 + 21 = 44$ kg.
}
\end{ex}

\begin{ex}
Một con cừu mẹ nặng 35 kg, cừu con nặng 12 kg. Cừu mẹ nặng hơn cừu con bao nhiêu ki-lô-gam?
\shortans{23}
\loigiai{
Nặng hơn: $35 - 12 = 23$ kg.
}
\end{ex}

\begin{ex}
Bác nông dân thu hoạch được 50 kg bí đỏ, bác đã bán đi 30 kg. Bác còn lại bao nhiêu ki-lô-gam bí đỏ?
\shortans{20}
\loigiai{
Còn lại: $50 - 30 = 20$ kg.
}
\end{ex}

\begin{ex}
Có 3 gói kẹo: gói thứ nhất nặng 2 kg, gói thứ hai nặng 3 kg, gói thứ ba nặng 4 kg. Ba gói nặng tất cả bao nhiêu ki-lô-gam?
\shortans{9}
\loigiai{
Cả ba gói nặng: $2 + 3 + 4 = 9$ kg.
}
\end{ex}
